\par\medskip\noindent
\SGAInlineRefTarget{sga3:xix:definition:1-6-1}\textbf{\textit{\SGARefLink{sga3:xix:definition:1:6:1}{Definition 1.6.1}.}}---
\textit{A \(k\)-group \(G\) is said to be \emph{reductive} if it is smooth,
affine, and connected, and if \(\operatorname{rad}(G)\) is a torus; that is,
if
\[
  \operatorname{rad}^{u}(G)=\{e\}.
\]}
\SGARefTarget{sga3:xix:definition:1:6:1}

\par\medskip\noindent
\textbf{\textit{\SGARefLink{sga3:xix:lemma:1:6:2}{Lemma 1.6.2}.}}---
\textit{Let \(G\) be a reductive \(k\)-group.
\begin{enumerate}
\item[\SGAInlineRefTarget{sga3:xix:definition:1:6:1:item:i}(i)]\SGAInlineRefTarget{sga3:xix:lemma:1:6:2:item:i} If \(Q\) is a torus of \(G\), then
  \(\underline{\operatorname{Centr}}_{G}(Q)\) is reductive.
\item[\SGAInlineRefTarget{sga3:xix:definition:1:6:1:item:ii}(ii)]\SGAInlineRefTarget{sga3:xix:lemma:1:6:2:item:ii} In particular, if \(T\) is a maximal torus of \(G\), then
  \(\underline{\operatorname{Centr}}_{G}(T)=T\).
\item[\SGAInlineRefTarget{sga3:xix:definition:1:6:1:item:iii}(iii)] The center of \(G\) is contained in every maximal torus and is
  therefore diagonalizable.
\item[\SGAInlineRefTarget{sga3:xix:definition:1:6:1:item:iv}(iv)]\SGAInlineRefTarget{sga3:xix:lemma:1:6:2:item:iv} The radical of \(G\) is the unique maximal torus of
  \(\underline{\operatorname{Centr}}(G)\).
\end{enumerate}}
\SGARefTarget{sga3:xix:lemma:1:6:2}

\SGAProof{Proof}
Indeed, \SGARefLink{sga3:xix:definition:1:6:1:item:i}{(i)} follows from (\SGARefLink{sga3:xix:eq:1:3:1}{1.3.1}), \SGARefLink{sga3:xix:definition:1:6:1:item:ii}{(ii)} from (\SGARefLink{sga3:xix:eq:1:3:2}{1.3.2}), and \SGARefLink{sga3:xix:definition:1:6:1:item:iii}{(iii)} follows from
\SGARefLink{sga3:xix:definition:1:6:1:item:ii}{(ii)}.  Finally, the maximal torus of
\(\underline{\operatorname{Centr}}(G)\), that is, the neutral component
\(\underline{\operatorname{Centr}}(G)^{0}\), is a smooth connected solvable
subgroup normal in \(G\), and hence is contained in
\(\operatorname{rad}(G)\).  Conversely, \(\operatorname{rad}(G)\) is a torus
normal in \(G\), and hence central (\emph{Bible}, §\SGARefLink{sga3:xix:paragraph:4-3}{4.3}, Corollary to
Proposition 2); this proves \SGARefLink{sga3:xix:definition:1:6:1:item:iv}{(iv)}.
\hfill\(\square\)

\par\medskip\noindent
\SGAInlineRefTarget{sga3:xix:remark:1-6-3}\textbf{\textit{\SGARefLink{sga3:xix:remark:1:6:3}{Remark 1.6.3}.}}---
\textit{If \(G\) is reductive and
\(\dim(G)\ne\operatorname{rgred}(G)\), then
\[
  \dim(G)-\operatorname{rgred}(G)\geq 2.
\]
Indeed, this difference is always even (cf.\ \SGARefLink{sga3:xix:paragraph:1-10}{1.10} below).}
\SGARefTarget{sga3:xix:remark:1:6:3}

\par\medskip\noindent\SGAInlineRefTarget{sga3:xix:paragraph:1-7}\textbf{1.7.}\quad
Let \(G\) be a smooth affine connected \(k\)-group, and let \(H\) be a
smooth connected normal subgroup.  Then
\begin{equation}
  \operatorname{rad}(H)
  =
  \bigl(\operatorname{rad}(G)\cap H\bigr)_{\mathrm{red}}^{0},
  \qquad
  \operatorname{rad}^{u}(H)
  =
  \bigl(\operatorname{rad}^{u}(G)\cap H\bigr)_{\mathrm{red}}^{0},
  \tag{1.7.1}
\end{equation}
as is seen immediately.  In particular, if \(G\) is reductive, then so is
\(H\).
\SGARefTarget{sga3:xix:equation:1:7:1}

Let \(f:G\longrightarrow G'\) be a surjective morphism of smooth affine
connected \(k\)-groups.  Then
\begin{equation}
  f\bigl(\operatorname{rad}^{u}(G)\bigr)
  =
  \operatorname{rad}^{u}(G').
  \tag{1.7.2}
\end{equation}
In particular, if \(G\) is reductive, then so is \(G'\).
\SGARefTarget{sga3:xix:equation:1:7:2}

Let us prove (\SGARefLink{sga3:xix:equation:1:7:2}{1.7.2}).  First,
\(f\bigl(\operatorname{rad}^{u}(G)\bigr)\) is contained in
\(\operatorname{rad}^{u}(G')\).  On introducing
\[
  H=
  \bigl(f^{-1}(\operatorname{rad}^{u}(G'))\bigr)_{\mathrm{red}}^{0},
\]
which contains \(\operatorname{rad}^{u}(G)\), we have
\(\operatorname{rad}^{u}(H)=\operatorname{rad}^{u}(G)\); thus we are reduced
to the case \(G=H\), that is, to the case in which \(G'\) is unipotent.
As \(T\) runs over the maximal tori of \(G\), the union of the groups
\(T\operatorname{rad}^{u}(G)\) is dense in \(G\).  Hence the union of the
groups \(f(T)\,f\bigl(\operatorname{rad}^{u}(G)\bigr)\) is dense in \(G'\).
But \(f(T)\) consists of semisimple elements, so \(f(T)=\{e\}\), since
\(G'\) is unipotent.  It follows that
\(f\bigl(\operatorname{rad}^{u}(G)\bigr)\) is dense in \(G'\).  Therefore,
by \emph{Bible}, §\SGARefLink{sga3:xix:paragraph:5:4}{5.4}, \SGARefLink{sga3:xix:lemma:1:4}{Lemma 4}, or [Ch05], §6.1, \SGARefLink{sga3:xix:lemma:3:1}{Lemma 1},\footnote{Editor's
note: the following argument has been supplied in detail.}
\(f\bigl(\operatorname{rad}^{u}(G)\bigr)\) is an open subgroup of \(G'\).
Since \(G'\) is connected, it follows that
\[
  f\bigl(\operatorname{rad}^{u}(G)\bigr)=G'.
\]

\noindent\textit{N.B.}\quad Under the same hypotheses one can prove that
\[
  f\bigl(\operatorname{rad}(G)\bigr)=\operatorname{rad}(G').
\]
